\guidepart{rose}{Fix it}{When things go wrong}
  {Symptom first, fix second. No essays.}
  {\partitem{\iLifebuoy}{Troubleshooting, by symptom}}
\label{part:fix}

% =========================================================================================
\chapter{Troubleshooting, by symptom}\label{ch:trouble}
\chapmeta{as long as it takes}{You'll find your symptom and the fix in one place}

\begin{tldr}
Find the heading that matches what you're seeing. Try the fixes in order. Most problems are the network, the
password, or the baud rate.
\end{tldr}

\section{Connecting}

\fix{``Authentication failed''}{%
\begin{itemize}
  \item Check the username, password, key path and the ticked \ui{Auth Methods}.
  \item If the key has a passphrase, check how the passphrase is stored (\ui{Passphrase Mode}).
  \item Two-factor server? Click \ui{Test} next to the TOTP secret and compare with your phone. Check your clock.
  \item Saved password not being used? The vault might be locked: click \ui{Vault locked} in the status bar.
\end{itemize}}

\fix{``Host key verification failed''}{%
The server's fingerprint changed since you last connected. That happens when a server is rebuilt, but it's also exactly
what an attack looks like. \textbf{Check with whoever runs the server.} If it's legitimate:
choose \ui{Accept \& Update} in the warning (or \ui{Accept Once} to defer the decision). Under a Strict or Accept-new
policy there's no prompt: \menu{Settings\then Known Hosts}, \ui{Remove} the entry and reconnect.}

\fix{It just says ``Connecting…'' forever}{%
\begin{itemize}
  \item Click \ui{Test} in the profile editor. If that fails too, it's the network, a firewall or the wrong port.
  \item Behind a corporate proxy? Set one up under \ui{Proxy} (SSH profiles).
  \item Going through a bastion? Check each jump host on its own first.
\end{itemize}}

\fix{The session keeps reconnecting}{%
Check the network and the credentials, then the profile's auto-reconnect settings. Reconnection backs off (1\,s, 2\,s,
4\,s… up to 30\,s) and stops after the maximum number of attempts.}

\section{Serial}

\fix{The port won't open}{%
\begin{itemize}
  \item Is the adapter plugged in? Hit refresh in the device list.
  \item \textbf{Permission denied} on Linux: add yourself to \code{dialout} (or \code{uucp}), then log out and back in.
  \item On a Mac, use \code{/dev/cu.*}, not \code{/dev/tty.*}.
  \item \textbf{Port is in use by another program}: close \code{screen}, \code{minicom}, the Arduino IDE serial monitor or
        whatever else has it open.
\end{itemize}}

\fix{Gibberish on screen}{The baud rate, parity or stop bits don't match the device. Check its manual.}

\section{Remote desktops}

\fix{VNC won't connect}{%
\begin{itemize}
  \item VNC ports are 5900 plus the display number: \code{:0} is 5900, \code{:1} is 5901.
  \item Check the password and how it's stored.
  \item Using an SSH tunnel? Check the tunnel's host, user, login method and key, and the host key policy.
  \item Connecting to a Mac? Turn on password access in macOS Screen Sharing settings (see p.\,\pageref{ch:vnc}).
\end{itemize}}

\fix{RDP won't connect}{%
\begin{itemize}
  \item Linux box running xrdp? Tick \ui{Allow servers without NLA}.
  \item ``Logon failed''? Check the domain. Kerberos users: the KDC must be reachable directly, not through the proxy.
  \item Certificate warning after a server rebuild? That's expected once; compare the fingerprint, then
        \ui{Accept \& Update}.
  \item Going through an RD Gateway? It needs a certificate from a public CA and Basic authentication over HTTPS.
\end{itemize}}

\section{Files}

\fix{Upload is greyed out or fails}{%
The folder is probably \proto[amber]{READ ONLY} for your user. Pick another, or check permissions on the server. Uploads
need an SSH session; other protocols can't do it.}

\fix{``NO READ'' on a folder}{You don't have permission to look inside. Try the parent folder, or ask the server's admin.}

\section{The terminal itself}

\fix{The terminal is the wrong size}{%
Resize the window or nudge a split divider; the terminal re-fits itself. Still wrong? Reconnect. Very old servers may
prefer a different \ui{Terminal Type} in the profile (try \code{vt100}).}

\fix{The scroll wheel moves the cursor instead of scrolling}{%
You're in a full-screen program (\code{less}, \code{vim}, \code{tmux}). That's how terminals work: the wheel sends arrow keys
there. Quit the program, or use its own scrolling (in \code{tmux}: \keys{Ctrl+B} then \kbd{[}).}

\fix{Emoji are in black and white}{Known limitation: colour emoji show as monochrome glyphs. Everything else renders
normally.}

\begin{quickwin}[Still stuck?]
Start Nexus Terminal with \code{--data-dir} pointing at an empty folder and try the same connection there. If it
works, something in your profile is off; duplicate the working one and compare.
\end{quickwin}

\section{Still stuck?}

\fix{Something odd that isn't on this list}{%
Look at \code{nexus.log} in your data folder (\menu{Settings\then Diagnostics} shows where): connections, reconnects,
host-key checks and errors, never passwords or what's on screen. Start Nexus Terminal with \code{RUST\_LOG=debug} for
more detail. Reporting a bug? \ui{Copy Diagnostics} and paste it in.}

